\documentclass[RNAAS]{aastex701}
% ---------------------------------------------------------------------------------------------
% RELEASE-DAY TEMPLATE. Every \dr{...} marks a number filled in after the Gaia DR4 analysis
% (2 December 2026). Every \fz{...} marks a number fixed at the catalogue freeze (27 November, 03:45 UTC; N_new is set when the SCoPe check of all 5,778 candidates finishes, before DR4).
% Remove the two macros' highlighting before submission. Limit: 1,500 words incl. references,
% abstract <= 150 words, one table, first person singular.
% ---------------------------------------------------------------------------------------------
\usepackage{xcolor}
\newcommand{\dr}[1]{\textcolor{red}{[DR4: #1]}}
\newcommand{\fz}[1]{\textcolor{blue}{[#1]}}

\begin{document}

\title{Variable stars hidden from Gaia DR3: a pre-registered ZTF census tested against Gaia DR4}

\author[0000-0001-8779-0700]{Deepnil Ray}
\affiliation{Ashoka University, Sonipat, Haryana 131029, India}
\email{deepnilray_ug2025@ashoka.edu.in}
\affiliation{The NoRCEL Institute, United Kingdom}

\begin{abstract}
Gaia DR3 published variability results for about 10.5 million sources and none for the rest. I searched
\fz{13.9 million} stars that DR3 left without any variability result, using Zwicky Transient Facility (ZTF) DR23 light
curves and a detector validated on sky tiles sealed before any data were opened (AUC 0.954; false-positive rate
0.23\%). The search yields \fz{5,778} periodic variables confirmed in independent light curves and stable season by
season. Before Gaia DR4 appeared, I deposited the catalogue and a dated prediction of which stars DR4 would flag.
DR4 flags \dr{X\%} of them as variable against \dr{Y\%} of matched controls, and agrees with my period for
\dr{Z\%}. \fz{N\_new} stars appear in no published catalogue, including the ZTF SCoPe classifications; these are new
variables. The sample measures directly which real variables Gaia's pipeline misses.
\end{abstract}

\keywords{Variable stars (1761) --- Sky surveys (1464) --- Astronomy data analysis (1858)}

\section{Motivation}
Gaia's sparse sampling, set by its scanning law, leaves real variables unflagged, but how many, and of which
periods and amplitudes, has not been measured; \citet{Everall2022} describe a selection function for Gaia's
variable samples as ``hugely valuable'' and note that building one requires knowing which sources Gaia would have
classified as variable had they been variable. A census of real variables among stars Gaia called quiet answers
that question directly. Testing such a census against the next release, with the prediction fixed in advance,
makes the answer verifiable.

\section{Data and method}
I selected Gaia DR3 stars with no \texttt{vari\_summary} row, not flagged as variable, and RUWE $<1.4$, at
$|b|\ge15^\circ$ and $\delta\ge-25^\circ$, and read their ZTF DR23 light curves from the IRSA HATS archive. A
seeded manifest assigned every HEALPix tile to training, validation or search before any download. A gradient-boosted
classifier on light-curve shape features (Lomb--Scargle power and false-alarm probability, amplitude over error,
reduced $\chi^2$, skewness, von Neumann ratio and an eclipse box-fit channel) was frozen before the validation tiles
were opened. On 356 Gaia-classified variables and 4,000 Gaia-quiet stars in those tiles it reached AUC 0.954
(lower 95\% bound 0.939; shuffled-label control 0.494). The full search rule has a false-positive rate of 0.23\%
(upper bound 0.43\%), itself an upper limit because some ``negatives'' are unrecognised variables. Every candidate
had to recover its period in at least two independent ZTF light curves, in at least half of the seasons tested,
away from sampling aliases, with no brighter Gaia star within $6''$. Only grades A and B are used.

\section{Results}
Table~\ref{tab:res} summarises the census and its test. Prior art was checked against VSX, SIMBAD, ALeRCE, 17
VizieR catalogues and the latest SCoPe release \citep{Healy2024}. SCoPe already flags most candidates as
variable, independently confirming that the detections are real; for \fz{11.8\%} SCoPe's period or class disagrees
with mine, and \fz{0.9\%} are absent from SCoPe altogether. The \fz{N\_new} stars absent from every catalogue are new
variable stars and are being reported to VSX.

On 2 December 2026, Gaia DR4 \citep{GaiaDR4} classified the frozen stars as follows: \dr{N\_VAR} VARIABLE,
\dr{N\_POS} POSSIBLE and \dr{N\_NON} NON\_VARIABLE, compared with \dr{control fractions} for the matched
control list. DR4 periods match mine (within 1\%, allowing factors of 0.5 and 2) for \dr{fraction}. Of the \fz{N\_new}
new variables, DR4 flags \dr{n}.

\begin{deluxetable}{lc}
\tablecaption{Census of Gaia-DR3-quiet variables and its test against Gaia DR4.\label{tab:res}}
\tablehead{\colhead{Quantity} & \colhead{Value}}
\startdata
Gaia-DR3-quiet stars searched in ZTF & \fz{13.9 million} \\
Validated A/B periodic variables & \fz{5,778} \\
\quad already in SCoPe with matching class or period & \fz{87.3\%} \\
\quad in SCoPe with different period or class & \fz{11.8\%} \\
\quad in no catalogue (new) & \fz{N\_new} \\
Flagged VARIABLE or POSSIBLE by Gaia DR4 & \dr{X\%} \\
Control stars flagged by Gaia DR4 & \dr{Y\%} \\
DR4 period agrees with ZTF period & \dr{Z\%} \\
\enddata
\end{deluxetable}

\section{Discussion}
A census fixed before the data release and then confirmed by it is the strongest test available that the
detections are real. The \dr{X\%} recovery shows that DR4, with twice the baseline of DR3, recovers most of the
variables DR3 missed; the remaining \dr{100-X\%} define where Gaia's variability detection still fails, which is
the input a variability selection function needs. A full selection function in period, amplitude, magnitude and
scanning-law coverage will be presented in a forthcoming paper.

\section*{Data availability}
The frozen catalogue, the pre-registered prediction and the code are archived on Zenodo (\fz{DOI}, deposited
27 November 2026) and at \url{https://github.com/deepnilray/SearchChev}.

\begin{acknowledgments}
This work used ZTF data (IRSA) and data from the ESA mission Gaia, processed by the Gaia DPAC. Software development
and analysis were assisted by an AI system (Claude, Anthropic); all results were checked by the author.
\end{acknowledgments}

\begin{thebibliography}{}
\bibitem[Everall \& Boubert(2022)]{Everall2022} Everall, A., \& Boubert, D. 2022, MNRAS, 509, 6205
\bibitem[Gaia Collaboration(2026)]{GaiaDR4} Gaia Collaboration 2026, A\&A, \dr{reference}
\bibitem[Healy et al.(2024)]{Healy2024} Healy, B. F., et al. 2024, ApJS, 272, 14
\end{thebibliography}

\end{document}
